\documentclass[12pt, a0paper]{article}

% ===== 页面设置 =====
\usepackage[
    paperwidth=84.1cm,
    paperheight=135cm,
    margin=1.5cm,
    columnsep=1.8cm
]{geometry}

% ===== 中文 =====
\usepackage{ctex}

% ===== 双栏 =====
\usepackage{multicol}

% ===== 表格颜色 =====
\usepackage[table]{xcolor}
\usepackage{colortbl}
\usepackage{multirow}

% ===== 标题格式 =====
\usepackage{titlesec}
\usepackage{fancyhdr}
\usepackage{tikz}
\usepackage{enumitem}
\usepackage{amsmath}

% ===== 颜色 =====
\definecolor{myblue}{RGB}{25,55,109}
\definecolor{myorange}{RGB}{200,100,20}
\definecolor{mygreen}{RGB}{30,120,70}
\definecolor{myred}{RGB}{180,40,40}
\definecolor{mygray}{RGB}{130,130,130}
\definecolor{mypurple}{RGB}{100,40,160}
\definecolor{lightblue}{RGB}{235,245,255}

% ===== 字体缩放（A0 需要大字体）=====
\renewcommand{\normalsize}{\fontsize{22}{30}\selectfont}
\renewcommand{\large}{\fontsize{28}{36}\selectfont}
\renewcommand{\Large}{\fontsize{34}{42}\selectfont}
\renewcommand{\LARGE}{\fontsize{40}{48}\selectfont}
\renewcommand{\huge}{\fontsize{48}{56}\selectfont}
\renewcommand{\Huge}{\fontsize{58}{66}\selectfont}
\renewcommand{\small}{\fontsize{18}{24}\selectfont}
\renewcommand{\footnotesize}{\fontsize{15}{20}\selectfont}

% ===== 收紧间距 =====
\setlength{\parskip}{0.3em}
\setlength{\parindent}{0pt}
\setlength{\multicolsep}{0.5em}
\raggedbottom

% ===== 节标题格式 =====
\titleformat{\section}[hang]
    {\huge\bfseries\color{myblue}}
    {\thesection}
    {1em}
    {}
    [\vspace{-1mm}{\color{myblue}\rule{\columnwidth}{1.5pt}}]
\titlespacing{\section}{0pt}{0.5em}{0.3em}

\titleformat{\subsection}[hang]
    {\Large\bfseries\color{myblue}}
    {\thesubsection}
    {1em}
    {}
\titlespacing{\subsection}{0pt}{0.3em}{0.1em}

% ===== 列表间距 =====
\setlist{nosep, leftmargin=3em, itemsep=0.5em}

% ===== 双栏脚注线 =====
\pagestyle{fancy}
\renewcommand{\headrulewidth}{0pt}
\renewcommand{\footrulewidth}{0pt}
\fancyfoot{}

% ===== 占位图命令 =====
\newcommand{\placeholder}[2]{%
    \begin{center}
    \colorbox{mygray!12}{\parbox{0.9\columnwidth}{%
        \centering
        \vspace{1.5em}
        {\Large\bfseries\color{mygray} [ 占位图： #1 ]}
        \vspace{1em}

        {\small\color{mygray} #2}
        \vspace{1.5em}
    }}
    \end{center}
}

% ===== 高亮框 =====
\newcommand{\highlight}[1]{%
    \colorbox{lightblue}{\parbox{\columnwidth}{\centering\bfseries #1}}
}

\begin{document}

% =====================
% 标题区（跨双栏）
% =====================
\begin{center}
    {\fontsize{64}{76}\selectfont\bfseries\color{myblue}
        从任务专用模型到基础模型：\\[5mm]
        EEG 解码的泛化能力演进与可解释性挑战\\[12mm]
    }

    {\fontsize{32}{40}\selectfont
        第一作者 (SUSTech ID) \hspace{2em} 第二作者 (ID) \hspace{2em} 第三作者 (ID)\\[5mm]
    }

    {\fontsize{24}{32}\selectfont\color{mygray}
        南方科技大学生物医学工程系 \hspace{3em}
        BMEB215 机器学习与医学工程应用 · 2026 Spring
    }

    \vspace{8mm}
    {\color{myblue}\hrule height 3pt}
    \vspace{8mm}
\end{center}

% =====================
% 双栏正文
% =====================
\begin{multicols}{2}

% ================================================================
% 1. 引言
% ================================================================
\section{引言：为什么需要更强的 EEG 模型？}

EEG 在脑机接口、临床诊断、睡眠分期、情绪识别等领域应用广泛，但其解码长期面临三个核心困难：

\textbf{标注昂贵。}临床诊断、睡眠分期、运动想象等任务需要专家标注或严格实验设计，标签成本高，难以大规模扩充。

\textbf{被试差异大。}EEG 信号受被试、session、设备、导联、采样率和预处理流程的影响明显，同一任务在不同条件下的分布偏移严重。

\textbf{泛化能力差。}传统任务专用模型通常在特定数据集上从零训练或强监督训练，换一个任务往往需要重新设计、重新调参，跨任务复用困难。

面对这些挑战，EEG 解码模型经历了从紧凑 CNN、到 CNN+Transformer 混合架构、再到大规模预训练基础模型的演进。本文沿着这条主线展开，并追问：\textbf{模型越来越强之后，我们失去了什么？}

% ================================================================
% 2. 模型演进
% ================================================================
\section{模型演进：Foundation Model 是如何出现的？}

\subsection{第一阶段：EEGNet（2018）—— 紧凑 CNN 专家模型}

EEGNet 不是简单的“CNN 阶段”，而是深度学习进入 EEG 解码领域的代表性工作。它用 Temporal Convolution 学习频率相关的时间滤波器，Depthwise Spatial Convolution 学习通道空间模式，再加 Separable Convolution 降参，形成紧凑而领域感知的结构。每个下游任务单独监督训练，在单任务上表现强劲，但局限也很明显——感受野偏局部，跨任务复用弱，仍是任务专用训练范式。

\subsection{第二阶段：EEG Conformer（2023）—— CNN + Transformer}

EEG Conformer 在卷积局部特征提取之后引入 Transformer Encoder，用 Self-Attention 建模 Patch 之间的全局依赖。这一步的意义不仅是结构升级：自注意力支持并行计算，参数可扩展，使模型开始具备适配大规模预训练的结构基础。在本项目中，我们还进行了“去掉 Transformer Block”的消融实验，以检验 Attention 对单任务 EEG 解码的贡献。

\subsection{第三阶段：CBraMod（2025）—— EEG Foundation Model}

CBraMod 用 Criss-Cross 注意力机制和大规模 EEG 预训练，实现了从“每个任务训练专家模型”到“预训练通用表征再适配下游任务”的范式转变。其 Backbone 输出 Channel-Patch 表征 $B \times C \times S \times D$，下游通过 Linear Probing（LP，冻结 Backbone）或 Full Fine-Tuning（FT）适配不同任务。LP 衡量预训练表征本身的可读性，FT 衡量最终下游适配能力。

\vspace{2em}
\highlight{核心结论：模型能力不断增强，泛化能力不断提高。从 CNN 局部特征，到 Transformer 全局建模，再到预训练通用表征。}

\vspace{2em}

\placeholder{模型架构对比}{%
    EEGNet / EEG Conformer / CBraMod 结构示意图。\\
    展示三类模型的输入、骨干网络、输出头结构差异。\\
    \small (等待绘制或队友提供)}

% ================================================================
% 3. 可解释性挑战
% ================================================================
\section{可解释性挑战：模型变强后，我们失去了什么？}

模型能力的提升伴随着表征的不断抽象化。

\textbf{传统 EEG 分析}依赖显式生理特征：Alpha、Beta、Theta、Delta 频段的能量分布，时域统计量（均值、方差、峰度），复杂度特征（样本熵、Lempel-Ziv 复杂度）等。这些特征直接对应可解释的生理量，研究者可以明确说出模型在用什么信息做决策。

\textbf{深度特征}虽然可以部分可视化（如 CNN 卷积核的频带偏好、Transformer Attention Map 的时间-通道关注模式），但语义已经不再明确——一个卷积核可能混合了多个频带的信息，一个 Attention Head 可能关注的是“某种统计模式”而非特定的生理事件。

\textbf{Foundation Embeddings}将这一问题推到了极致。CBraMod 等基座模型输出的高维隐空间向量泛化能力极强，但单看一个 Embedding 维度几乎无法对应任何经典的 EEG 描述符。基座模型的表征虽然泛化能力强，但可解释性正成为一个不容忽视的问题。

\begin{center}
\colorbox{myorange!10}{\parbox{0.95\columnwidth}{\centering\bfseries
    表征越来越抽象 \quad $\Rightarrow$ \quad 解释难度增大}}
\end{center}

\vspace{1em}

\textbf{措辞注意：}本海报使用“Interpretability Challenge”（可解释性挑战），而非“可解释性下降”。我们不需要证明它一定下降了——只需指出，模型学到的表征越来越难以被人类理解，这正在成为下一阶段的关键问题。

\placeholder{可解释性挑战示意图}{%
    传统 EEG 特征（Alpha/Beta/Theta 频段 · 可解释生理量）\\
    $\downarrow$\\
    深度特征（CNN 卷积核 · Attention Map · 语义不明确）\\
    $\downarrow$\\
    Foundation Embeddings（高维隐空间向量 · 泛化强但难以理解）\\
    $\downarrow$\\
    抽象程度递增，解释难度增大。\\
    \small (等待绘制)}

% ================================================================
% 4. 对比分析
% ================================================================
\section{对比总结：一张表看完全文}

\begin{center}
{\small\itshape\color{mygray} 评委即使只看这张表也能理解核心观点}

\vspace{1em}

\begin{tabular}{|l|c|c|c|}
    \hline
    \textbf{模型} & \textbf{核心思想} & \textbf{泛化能力} & \textbf{可解释性} \\
    \hline
    EEGNet        & CNN               & \textcolor{myred}{低}    & \textcolor{mygreen}{高}     \\
    \hline
    EEG Conformer & CNN + Transformer & \textcolor{myorange}{中}  & \textcolor{myorange}{中}    \\
    \hline
    CBraMod       & 大规模预训练      & \textcolor{mygreen}{高}   & \textcolor{myred}{较低}     \\
    \hline
    \rowcolor{mygreen!10}
    \textbf{+ NeuroGate} & \textbf{+ 神经科学先验} & \textcolor{mygreen}{\textbf{高}} & \textcolor{myblue}{\textbf{预期更高}} \\
    \hline
\end{tabular}
\end{center}

\vspace{1em}

模型越强 $\rightarrow$ 泛化越好 $\rightarrow$ 但可解释性越难。这张表是整张海报的压缩版。

\placeholder{CBraMod 复现结果}{%
    CBraMod 原论文任务复现主表。\\
    包含 BCIC2A、FACED、ISRUC-S1、Physionet-MI、SEED-V、SHU-MI、MDD、TUEV 等数据集。\\
    每项报告 3-seed avg bAcc、best bAcc、论文参考值。\\
    \small (等待队友复现结果填入)}

% ================================================================
% 5. 未来方向：NeuroGate
% ================================================================
\section{未来方向：NeuroGate}

\subsection{为什么需要 NeuroGate？}

基座大模型的一个重要问题是可解释性不足。CBraMod 的 Embedding 虽然泛化能力强，但高维隐空间不直接暴露频带能量、时域统计、复杂度等可解释生理量。NeuroGate 的思路是用显式 EEG 知识补充 Foundation Embedding：提取人工神经先验特征，通过 Gated Fusion 与 Backbone Embedding 融合。

\subsection{融合方式}

对每个 EEG Segment 提取人工神经先验特征向量 $m$（时域统计、Band Power、频谱形状、复杂度特征等）。将 CBraMod 输出的 Backbone Embedding 记为 $h$，用一个小的投影网络 $\phi$ 把 $m$ 映射到与 $h$ 可融合的空间，通过门控权重 $g$ 控制模型在多大程度上接受人工先验：

\[
z = h + g \cdot \phi(m)
\]

这样做的目的是让人工特征补充 Backbone Embedding，而非替代它。门控机制允许模型自主决定接受多少先验——在人工先验有用的任务上接受更多，在先验噪声大的任务上适度忽略。

\subsection{实验证据}

LP（冻结 Backbone）下，NeuroGate 在 HMC 睡眠分期（$+0.0773$）、ISRUC-S1（$+0.0972$）等依赖频谱节律结构的任务上提升明显，说明人工先验能补充 Frozen Embedding 中不易被线性头读出的信息。FT 下提升较小但总体为正。

\begin{center}
\colorbox{mypurple!8}{\parbox{0.95\columnwidth}{\centering\bfseries
    NeuroGate represents one possible direction.}}
\end{center}

\vspace{1em}

{\small
\textcolor{myred}{[不写]} NeuroGate 解决了可解释性问题。\hspace{2em}
\textcolor{myblue}{[而写]} NeuroGate represents one possible direction.
}

\placeholder{NeuroGate 结构示意图}{%
    CBraMod Backbone \quad + \quad 神经科学先验（频带能量 · 时域统计 · 复杂度）\\
    $\downarrow$\\
    Gated Fusion：$z = h + g \cdot \phi(m)$\\
    $\downarrow$\\
    增强表征（泛化能力 + 可解释性）\\
    \small (等待绘制)}

\placeholder{NeuroGate 实验结果}{%
    HMC / ISRUC-S1 / MDD / Physionet-MI 等数据集上的 LP 对比。\\
    CBraMod LP vs NeuroGate LP 柱状图或表格。\\
    \small (等待队友复现结果)}

% ================================================================
% 结论
% ================================================================
\section{如果评委只记住一句话}

\begin{center}
\colorbox{lightblue}{\parbox{0.95\columnwidth}{\centering\LARGE\bfseries
    EEG Foundation Models 正在解决泛化问题，\\
    而可解释性正在成为新的瓶颈。\\[5mm]
    \Large
    NeuroGate 尝试将神经科学知识重新引入 Foundation Models\\
    —— 其中一个可能的方向。}}
\end{center}

\vspace{2em}

% ================================================================
% 参考文献
% ================================================================
\section*{参考文献}

{\footnotesize
\begin{enumerate}[leftmargin=2em]
    \item Lawhern et al. \textbf{EEGNet: A Compact Convolutional Neural Network for EEG-based Brain-Computer Interfaces}. \textit{Journal of Neural Engineering}, 2018.
    \item Song et al. \textbf{EEG Conformer: Convolutional Transformer for EEG Decoding and Visualization}. \textit{IEEE Transactions on Neural Systems and Rehabilitation Engineering}, 2023.
    \item Wang et al. \textbf{CBraMod: A Criss-Cross Brain Foundation Model for EEG Decoding}. \textit{ICLR}, 2025.
    \item NeuroGate / NeuroKE 内部参考：neural-prior gated fusion experiments under LP and FT.
    \item 本地复现记录：\texttt{cbramod\_reproduction\_results\_summary.md}
\end{enumerate}
}

\end{multicols}

% =====================
% 底部线
% =====================
\vspace{5mm}
{\color{myblue}\hrule height 2pt}
\vspace{3mm}
\begin{center}
{\footnotesize\color{mygray}
    BMEB215 · 机器学习与医学工程应用 · 2026 Spring \hspace{3em}
    南方科技大学 · 生物医学工程系
}
\end{center}

\end{document}
